\documentclass[11pt]{article}

\usepackage[margin=1in]{geometry}
\usepackage{amsmath,amssymb}
\usepackage{booktabs}
\usepackage{array}
\usepackage{enumitem}
\usepackage[strings]{underscore}
\usepackage[hidelinks]{hyperref}

\setlist[itemize]{leftmargin=*}
\setlist[enumerate]{leftmargin=*}

\newcommand{\OffCEM}{\textsc{OffCEM}}
\newcommand{\DR}{\textsc{DR}}
\newcommand{\DM}{\textsc{DM}}
\newcommand{\IPS}{\textsc{IPS}}
\newcommand{\MIPS}{\textsc{MIPS}}
\newcommand{\relMSE}{\mathrm{relMSE}}
\newcommand{\dm}{\mathrm{dm}}

\title{Does OffCEM Need Oracle Cluster Labels?\\Evidence from Existing Experiments}
\date{}

\begin{document}
\maketitle

\section{Question}

The original \OffCEM{} synthetic benchmark constructs rewards from an action
partition. The practical question is whether \OffCEM{} works because it has
access to this reward-defining partition, or whether it can discover a useful
partition from observed action features.

The existing experiments support the following conclusion:

\begin{quote}
\OffCEM{} does not require a file of oracle labels as an input in principle, but
it does require the estimation partition to align closely with the
reward-relevant partition geometry. In the default benchmark, that geometry is
hidden and not recoverable by ordinary feature clustering. The feature-bucket
positive result is therefore best interpreted as a known-rule or oracle-like
control, not as evidence of practical cluster discovery.
\end{quote}

\section{Definitions}

Let $\phi_{\mathrm{gen}}$ be the reward-generating partition and
$\phi_{\mathrm{est}}$ be the partition used by \OffCEM{} during estimation. The
synthetic reward has the form
\[
q(x,a) = g(x,\phi_{\mathrm{gen}}(a)) + h_{\phi_{\mathrm{gen}}(a)}(x,a).
\]
Changing $\phi_{\mathrm{est}}$ changes the within-cluster residual target that
the two-stage reward model must learn. Thus the key issue is not whether a
method uses features, but whether its partition matches the reward residual
geometry.

We use the within-cluster demeaned error
\[
\dm(f;\phi) =
\mathbb{E}_x\left[
  \frac{1}{|\mathcal{A}|}
  \sum_a
  \left(
    f(x,a)-q(x,a)
    -
    \frac{1}{|\mathcal{A}_{\phi(a)}|}
    \sum_{b:\phi(b)=\phi(a)}(f(x,b)-q(x,b))
  \right)^2
\right]
\]
as the local-correctness diagnostic. This removes cluster-level offsets and
isolates whether the reward model preserves within-cluster reward contrasts.

\section{Evidence from the Default Benchmark}

In the default benchmark, rewards are generated using the paper's stochastic
``original'' partition. That partition is only weakly feature-dependent and has
near-zero agreement with standard feature clusterings. The local-correctness
results show that the exact generating partition is much better than any
mismatched alternative.

\begin{table}[h]
\centering
\caption{Local-correctness ratios on the default benchmark. Values are
$\dm_{\mathrm{2s}}/\dm_{\mathrm{1s}}$; smaller is better. Source:
\texttt{local_correctness_results/local_correctness_findings.md}.}
\begin{tabular}{lrrrr}
\toprule
Estimation partition & $K=10$ & $K=50$ & $K=100$ & $K=200$ \\
\midrule
matched & 0.61 & 0.41 & 0.20 & 0.14 \\
original rerun & 0.93 & 0.72 & 0.46 & 0.39 \\
random & 0.97 & 0.75 & 0.46 & 0.38 \\
shuffled partition & 0.94 & 0.68 & 0.46 & 0.36 \\
\bottomrule
\end{tabular}
\end{table}

At $K=50$, the exact matched partition reduces within-cluster error by roughly
60\%, while rerunning the paper's original stochastic mechanism behaves much
closer to random. This matters because the original mechanism redraws a new
partition; it is not the same as reusing the reward-generating labels.

The temperature sweep reinforces this point. Lowering the original mechanism's
softmax temperature changes cluster-size balance, but it does not make the
partition recoverable by K-means. The adjusted Rand index between the generating
partition and K-means remains approximately zero across temperatures, and
K-means is worse than random in the local-correctness ratio.

\section{Corruption Evidence}

The corruption sweep starts from the exact generating partition and randomly
reassigns a fraction $p$ of actions. This gives an interpretable alignment knob:
$p=0$ is the matched partition, and $p=1$ is close to a random partition.

\begin{table}[h]
\centering
\caption{Default-benchmark corruption sweep for local correctness at $K=50$.
Values are $\dm_{\mathrm{2s}}/\dm_{\mathrm{1s}}$. Source:
\texttt{local_correctness_results/corrupt_p_*/corrupt/nc50.json}.}
\begin{tabular}{rrrr}
\toprule
Corruption $p$ & $\dm_{\mathrm{2s}}$ & $\dm_{\mathrm{1s}}$ & ratio \\
\midrule
0.00 & 8.01 & 18.39 & 0.44 \\
0.10 & 9.50 & 19.39 & 0.49 \\
0.25 & 11.98 & 20.67 & 0.58 \\
0.50 & 14.09 & 22.26 & 0.63 \\
0.75 & 15.90 & 23.24 & 0.68 \\
1.00 & 16.76 & 23.54 & 0.71 \\
\bottomrule
\end{tabular}
\end{table}

This shows a smooth degradation rather than a cliff: as the estimation partition
moves away from the reward-generating partition, the local-correctness benefit
shrinks continuously. The JSON files for this sweep store estimator point values
but not the true policy value, so they do not by themselves contain relMSE.

The join experiments nevertheless show that this diagnostic is consequential.
Across the full 36-cell join grid, $\dm$ ranks \OffCEM{} error within each fixed
problem: mean Spearman correlation is about $+0.80$ for $\dm$ versus relative
bias squared, and the lowest-$\dm$ clustering is the lowest-bias clustering in
36/36 cells.

\section{Baseline Crossover}

The join results connect partition quality to actual OPE error. In the default
benchmark, well-aligned partitions beat \DR{}, while badly aligned partitions do
not.

\begin{table}[h]
\centering
\caption{Default-generation join results pooled over 36 cells. Source:
\texttt{join_results/cell_*.json} and the summary in
\texttt{local_correctness_results/local_correctness_findings.md}.}
\begin{tabular}{lrr}
\toprule
\OffCEM{} estimation partition & Mean \relMSE{} & Cells beating \DR{} \\
\midrule
matched & 1.32 & 36/36 \\
corrupt $p=0.25$ & 1.78 & 36/36 \\
corrupt $p=0.75$ & 2.41 & 16/36 \\
random / K-means / shuffled & 2.50--2.56 & 6--14/36 \\
\DR{} baseline & 2.41 & -- \\
\bottomrule
\end{tabular}
\end{table}

Thus poor clustering does not merely worsen an internal diagnostic. It can make
\OffCEM{} worse than a standard baseline. This supports the practical statement
that \OffCEM{} is useful only when the clustering is aligned enough.

\section{Feature-Bucket Result: Positive Control, Not Discovery}

The feature-coherent benchmark changes the generator so the reward partition is
a deterministic feature-bucket rule. This rule projects action features onto the
first principal component and slices the ordered actions into equal-frequency
buckets. Estimation with the same feature-bucket rule exactly reconstructs the
generating partition.

\begin{table}[h]
\centering
\caption{Feature-bucket-generated reward, $K=50$. Source:
\texttt{local_correctness_results/feature_coherent_fb/*/nc50.json}.}
\begin{tabular}{lrrr}
\toprule
Estimation partition & $\dm_{\mathrm{2s}}$ & $\dm_{\mathrm{1s}}$ & ratio \\
\midrule
matched oracle & 7.21 & 18.06 & 0.40 \\
feature bucket & 7.00 & 17.96 & 0.39 \\
K-means & 17.51 & 20.74 & 0.85 \\
random & 16.95 & 22.77 & 0.74 \\
\bottomrule
\end{tabular}
\end{table}

At full join scale, feature bucket also ties the matched oracle on estimator
error.

\begin{table}[h]
\centering
\caption{Feature-bucket-generation join results pooled over 36 cells. Source:
\texttt{join_results_fc/cell_*.json}.}
\begin{tabular}{lrr}
\toprule
\OffCEM{} estimation partition & Mean \relMSE{} & Cells beating \DR{} \\
\midrule
matched oracle & 1.34 & 36/36 \\
feature bucket & 1.34 & 36/36 \\
corrupt $p=0.25$ & 1.80 & 36/36 \\
corrupt $p=0.75$ & 2.46 & 7/36 \\
K-means & 2.49 & 8/36 \\
random & 2.57 & 0/36 \\
\DR{} baseline & 2.36 & -- \\
\bottomrule
\end{tabular}
\end{table}

This is an important positive control: literal label access is not necessary if
the correct partition rule is known and reproducible from observed features.
However, it is not evidence that a practitioner can generally discover the
right partition. The success comes from using the same deterministic rule during
generation and estimation. K-means uses the same action features but fails
because it captures a different feature geometry.

For that reason, the feature-bucket result should be described as a known-rule
or oracle-like structural control, not as an oracle-free discovery result.

\section{Reverse-Mismatch Evidence}

The existing imprecise-clustering pilot already contains the reverse mismatch
that would otherwise require a new experiment: generate rewards using
feature-bucket clusters, then estimate using the paper's original mechanism.

\begin{table}[h]
\centering
\caption{Reverse mismatch: feature-bucket generation with alternative estimation
partitions. Source: \texttt{imprecise_results/direction1/method_summary.csv}.
This is a 5-seed pilot, so it is supporting evidence rather than the main
significance-tested result.}
\begin{tabular}{lrrr}
\toprule
Generation & Estimation partition & \relMSE{} & ARI with generator \\
\midrule
feature bucket & matched & 0.581 & 1.000 \\
feature bucket & feature bucket & 0.581 & 1.000 \\
feature bucket & K-means & 0.942 & 0.016 \\
feature bucket & original rerun & 1.175 & $\approx 0$ \\
feature bucket & random & 1.336 & $\approx 0$ \\
\bottomrule
\end{tabular}
\end{table}

This directly shows that the paper's original clustering mechanism is not
intrinsically special. When it is not the reward-generating mechanism, it fails
like other mismatched partitions.

\section{What the Existing Results Can and Cannot Claim}

\subsection{Supported claims}

\begin{enumerate}
    \item On the default benchmark, \OffCEM{}'s advantage depends on using the
    reward-generating partition or a close approximation to it.
    \item Standard feature clustering does not recover that partition. K-means
    has near-zero ARI with the default generating partition and performs like a
    mismatched or random partition.
    \item Rerunning the paper's original stochastic clustering is not equivalent
    to oracle access. It redraws a different partition and behaves much closer
    to random than to matched.
    \item The feature-bucket benchmark proves that literal label access is not
    logically required, but only when the correct deterministic feature rule is
    known and exactly reproducible.
    \item The paper's original mechanism also fails in reverse mismatch:
    feature-bucket generation plus original estimation performs poorly relative
    to matched or feature-bucket estimation.
\end{enumerate}

\subsection{Claims that are not supported}

\begin{enumerate}
    \item The results do not show that ordinary unsupervised clustering can
    discover the reward-relevant partition.
    \item The feature-bucket result should not be summarized as ``\OffCEM{}
    works without oracle information'' without qualification. It is better
    summarized as ``\OffCEM{} works when the correct partition rule is known from
    observed structure.''
    \item The current dense corruption sweep does not contain relMSE because the
    local-correctness JSONs do not save the true policy value. Dense relMSE over
    $p \in \{0,.1,.25,.5,.75,1\}$ would be a useful final figure, but it is not
    required for the core oracle-label conclusion.
\end{enumerate}

\section{Final Conclusion}

The most defensible conclusion is:

\begin{quote}
\OffCEM{} is not a general clustering-discovery method. Its empirical advantage
depends on the estimation partition being aligned with the reward-relevant
partition. In the original benchmark, that partition is effectively oracle
information because it is hidden in the synthetic reward construction and is
not recovered by feature clustering. The feature-bucket benchmark shows an
escape hatch: if the reward-relevant partition is generated by a known,
deterministic feature rule, then the same rule can match oracle performance
without reading labels. But this is a known-structure setting, not evidence that
the correct clusters can be discovered automatically.
\end{quote}

\end{document}
